\documentclass[11pt,a4paper]{article}
\usepackage[margin=25mm]{geometry}
\usepackage[T1]{fontenc}
\usepackage{lmodern}
\usepackage{amsmath,amssymb,amsthm,mathtools}
\usepackage{booktabs,array,graphicx}
\usepackage[numbers,sort&compress]{natbib}
\usepackage{microtype}
\usepackage[colorlinks=true,linkcolor=blue,citecolor=blue,urlcolor=blue]{hyperref}
\usepackage{doi}
\hypersetup{pdftitle={A finite nonreciprocal counterexample to a noise-matched current large-deviation bound},pdfauthor={Xiaojian Zeng},pdfsubject={Coherent transport and current large deviations}}
\newcommand{\kb}{k_{\mathrm B}}
\newcommand{\ee}{\varepsilon}
\newcommand{\ii}{\mathrm i}
\newcommand{\dd}{\mathrm d}
\newcommand{\Tr}{\operatorname{Tr}}
\newcommand{\asinh}{\operatorname{arsinh}}
\newcommand{\atanh}{\operatorname{artanh}}
\newcommand{\KL}{D}
\newcommand{\ellambda}{\lambda_{\mathrm q}}
\newtheorem{proposition}{Proposition}
\newtheorem{lemma}{Lemma}
\title{A finite nonreciprocal counterexample to a noise-matched\newline current large-deviation bound}
\author{Xiaojian Zeng\\[3pt]
\small PBC School of Finance, Tsinghua University, 100190\\
\small\texttt{xiaojianzeng@foxmail.com}}
\date{}
\begin{document}
\maketitle
\begin{abstract}
An established thermodynamic bound on the variance of a particle current in a
coherent conductor with broken reciprocity suggests a corresponding extension
of the reciprocal large-deviation bound. We formulate the proposed extension
by matching its curvature to that variance inequality and give a
computer-assisted counterexample. The conductor is a stationary Hermitian
387-dot sample attached to three normal tight-binding leads and three fixed
Fermi reservoirs at a common positive temperature. Its contact self-energies,
finite filter width and full propagating band are retained exactly. At the
interior reversal current $j=-J$, a complete interval integration and a
fixed-field Fenchel argument establish a rate-function violation greater than
$0.019\ee/h$, where $\ee=\kb T_{\mathrm{ref}}$. A conservative analytic bound
on the exact dissipation constant removes any reliance on its numerical
minimizer. The arithmetic verification includes an exhaustive check of every
Taylor remainder used in the fixed certificate, after an intermediate
magnitude-bound defect was identified. The result excludes this particular
noise-matched single-current shape; it leaves the reciprocal theorem, the
nonreciprocal variance bound and the alternative joint-current question intact.
\end{abstract}

\section{The proposed extension and its scope}

Mean dissipation constrains fluctuations in coherent fermionic transport, but
the constraint depends on the observable and on reciprocity. Brandner and
Saito proved a thermodynamic upper bound on the particle-current rate function
when the transmission coefficients are symmetric
\cite{BrandnerSaitoJPA2025}. Their separate variance inequality applies also
when this symmetry is broken \cite{BrandnerSaitoPRL2025}. In the perspectives
section of the final large-deviation paper, they propose inserting the
nonreciprocal dissipation factor into the full bound so as to recover the
variance inequality. They also raise a distinct possibility involving the
joint rate function of all independent currents.

The first suggestion is qualitative. We therefore state explicitly the
curvature-matching interpretation tested here. Let $J>0$ be an oriented
single-terminal mean particle current and let $\sigma>0$ be the actual total
reservoir entropy-production rate. Write
\begin{align}
 R_\psi&=\coth\!\left(\frac{\sigma}{2\psi\kb J}\right),\qquad c_\psi=JR_\psi,
 \label{eq:R}\\
 B_\psi(j)&=(c_\psi+j)\log\frac{c_\psi+j}{c_\psi+J}
 +(c_\psi-j)\log\frac{c_\psi-j}{c_\psi-J}.
 \label{eq:B}
\end{align}
The assertion is
\begin{equation}
 I(j)\le B_{\psi_0}(j),\qquad -c_{\psi_0}<j<c_{\psi_0},
 \label{eq:assertion}
\end{equation}
for every stationary elastic noninteracting normal coherent conductor, every
finite number of leads and channels, and all fixed finite reservoir chemical
potentials and positive temperatures for which the steady integrals exist.
No transmission reciprocity is imposed. Equation~\eqref{eq:assertion} is our
precise formalization of the first suggestion, rather than a formula quoted
verbatim as an established result.

The constant is the exact infimum defined in the supplemental material of
Ref.~\cite{BrandnerSaitoPRL2025}:
\begin{equation}
 \psi_0=\inf_{\substack{0<x,y<1\\x\ne y}}
 \frac{\rho(x)-\rho(y)-\rho'(x)(x-y)}
 {(x-y)\asinh\!\bigl((x-y)/(x+y-x^2-y^2)\bigr)},
 \quad \rho(x)=-x\log x-(1-x)\log(1-x).
 \label{eq:psi}
\end{equation}
Continuous limiting values are understood. We do not replace this infimum
by a rounded decimal. If $C_2=\chi''(0)$ is the count-variance rate, the
published nonreciprocal noise inequality reads
\begin{equation}
 \frac{C_2}{J}\sinh\!\left(\frac{\sigma}{\psi_0\kb J}\right)\ge1.
 \label{eq:noise}
\end{equation}
Indeed, $B_\psi''(J)=\sinh[\sigma/(\psi\kb J)]/J$, so
\eqref{eq:R} gives the curvature matching in this ansatz. The reciprocal
theorem instead uses $R=\coth[\sigma/(4\kb J)]$; it is not being tested
outside its hypotheses.

Our counterexample uses all reservoir heats in $\sigma$. This matters because
the proved nonreciprocal bound involving a partial dissipation in
Ref.~\cite{BrandnerSaitoJPA2025} does not permit that quantity to be replaced
by total dissipation without the additional symmetry premise. Bounds on
arbitrary-current precision based on forward-backward event entropy
\cite{TesserSplettstoesser2026} use another input and do not supply
\eqref{eq:assertion}. A bounded primary-literature and version check on
13 October 2026 identified no later resolution of the precise assertion
above. This is not an exhaustive priority claim.

\section{Counting conventions and a decisive fixed-field test}

At energy $E$, $S_{ab}$ maps the incident channel of lead $b$ to the outgoing
channel of lead $a$, and $\mathcal T_{ab}=|S_{ab}|^2$ in our one-channel
sample. The extracted-particle current and the heat injected into the sample
are
\begin{align}
 J_a&=\frac1h\int_0^\infty\!\dd E\sum_b
       \mathcal T_{ab}(f_a-f_b),\label{eq:J}\\
 Q_a&=\frac1h\int_0^\infty\!\dd E\sum_b
       \mathcal T_{ab}(f_a-f_b)(E-\mu_a),\qquad
 \sigma=-\sum_a\frac{Q_a}{T_a},\label{eq:entropy}\\
 f_a(E)&=\frac{1}{1+\exp[(E-\mu_a)/(\kb T_a)]}.
\end{align}
Here $h$ is Planck's constant. For $F=\operatorname{diag}(f_a)$ and the
channel projector $V_0$, the long-time cumulant-generating rate and its
Legendre--Fenchel transform are
\begin{align}
 \chi_0(s)&=\frac1h\int_0^\infty\!\dd E\,
 \log\det[\mathbb I-F+F S^\dagger e^{-sV_0}S e^{sV_0}],\label{eq:CGF}\\
 I_0(j)&=\sup_{s\in\mathbb R}\{sj-\chi_0(s)\}.
 \label{eq:LF}
\end{align}
Thus $s$ is dimensionless, $J,\chi,I$ have units of inverse time and
$\sigma$ has units of entropy per time. There is no extra factor of two
from an electrical-noise spectral-density convention.

For one counted channel the energy-resolved count $X$ takes values
$-1,0,1$. Direct conditioning on the incoming occupation of that channel,
or the Slater-determinant scattering probabilities, gives
\begin{align}
 p_+&=f_0\sum_{b\ne0}\mathcal T_{0b}(1-f_b),&
 p_-&=(1-f_0)\sum_{b\ne0}\mathcal T_{0b}f_b,&
 p_0&=1-p_+-p_-,\label{eq:marginal}\\
 M(s)&=p_0+p_+e^s+p_-e^{-s}.\label{eq:M}
\end{align}
These probabilities reduce the determinant in \eqref{eq:CGF}; they are not
an independent stochastic replacement of the scattering model.

The reversal current gives a particularly economical test. For $J,\sigma>0$
and finite positive $\psi$, $c_\psi>J$ and direct substitution yields
\begin{equation}
 B_\psi(-J)=2J\log\frac{c_\psi+J}{c_\psi-J}
           =\frac{2\sigma}{\psi\kb}.
 \label{eq:reversalB}
\end{equation}
Appendix~\ref{app:psi} proves the exact conservative bound
$1/2\le\psi_0\le1$. Fenchel's inequality at the fixed field $s=-24$
therefore implies
\begin{equation}
 I_0(-J)-B_{\psi_0}(-J)
 \ge24J-\chi_0(-24)-\frac{4\sigma}{\kb}
 \equiv\Delta.
 \label{eq:Delta}
\end{equation}
A positive enclosure of $\Delta$, together with $J,\sigma>0$, proves a
strict rate-function violation at an interior current. No numerical saddle,
rounded dissipation constant or assumption about equality in
\eqref{eq:LF} is required.

\section{A finite stationary normal conductor}
\label{sec:model}

Fix any $T_{\mathrm{ref}}>0$ and use the energy unit
$\ee=\kb T_{\mathrm{ref}}$. Each of three identical normal
one-channel semi-infinite leads has onsite $2t=20000\ee$ and hopping
$-t=-10000\ee$. Its propagating band is $0<E<40000\ee$.
The retarded surface Green function is
\begin{equation}
 g_{\mathrm{lead}}(E)=
 \frac{E-2t-\ii\sqrt{E(4t-E)}}{2t^2},\qquad 0<E<4t,
 \label{eq:surface}
\end{equation}
with the retarded analytic continuation outside the band. It follows from
the semi-infinite surface recurrence and equivalently from the outgoing
Bloch solution $E=2t-2t\cos k$.

The sample consists of three identical $n=128$ site filter chains and three
core sites, hence $3n+3=387$ sites. Set
\begin{align}
 E_0&=5000\ee,&w&=0.02\ee,&\Gamma_c&=200000\ee,\nonumber\\
 \Gamma_f&=\frac{w}{4\sin[\pi/(2n)]},&
 h_\ell&=\frac{w}{4\sqrt{\sin[(2\ell+1)\pi/(2n)]
                         \sin[(2\ell-1)\pi/(2n)]}},
 \label{eq:filter}\\
 \lambda&=\sqrt{\Gamma_f\Gamma_c},&&1\le\ell\le n-1.\nonumber
\end{align}
The filter profile uses the $w/4$ convention of the quantum-dot construction
in Ref.~\cite{BrandnerSaitoPRL2025}. These equations, rather than any
approximate transmission-window description, define the sample.

Label the filter sites $(a,\ell)$ and the core sites $c_a$, with $a=0,1,2$.
The filter hoppings are the real $h_\ell$. The real hopping $\lambda$ joins
$(a,n)$ to $c_a$. All sample onsite energies are $E_0$, except that of the
outer site $(a,1)$, which is $E_0-\operatorname{Re}\Sigma(E_0)$. The real
lead contact is defined by
\begin{equation}
 V^2=\frac{2t^2\Gamma_f}{\sqrt{E_0(4t-E_0)}},\qquad
 \Sigma(E)=V^2 g_{\mathrm{lead}}(E).
 \label{eq:contact}
\end{equation}
On the three core sites the remaining hopping matrix is
\begin{equation}
 H_{\mathrm{core}}-E_0\mathbb I=\ii\Gamma_c(P-P^{\mathsf T}),
 \qquad P=\begin{pmatrix}0&0&1\\1&0&0\\0&1&0\end{pmatrix}.
 \label{eq:core}
\end{equation}
It is Hermitian and carries a nonzero magnetic flux. All hoppings, contacts
and onsite energies are fixed and time independent. There are no
interactions, anomalous pairing terms, drives, probes or internal baths.

For $0<E<4t$, let $\Gamma(E)=-\operatorname{Im}\Sigma(E)>0$ and
\begin{equation}
 G(E)=[E\mathbb I-H_{\mathrm{sample}}-\Sigma(E)P_{\mathrm{out}}]^{-1},
 \qquad S_{ab}(E)=\delta_{ab}-2\ii\Gamma(E)G_{(a,1),(b,1)}(E),
 \label{eq:S}
\end{equation}
where $P_{\mathrm{out}}$ projects onto the three outer filter sites.
The resolvent identity
$G-G^\dagger=-2\ii\Gamma G P_{\mathrm{out}}G^\dagger$ proves unitarity
of this flux-normalized scattering matrix. Outside the propagating band the
leads have no open channels, so there is no reservoir-current contribution.
Discrete bound sample modes carry no extensive steady reservoir current in
this stationary noninteracting model. The lead endpoints can be handled by
their limiting scattering values and have zero integration measure.

The three reservoirs remain at the common temperature $T_{\mathrm{ref}}$
and at the fixed chemical potentials
\begin{equation}
 (\mu_0,\mu_1,\mu_2)=(5001,4997,4976)\ee.
 \label{eq:reservoirs}
\end{equation}
No energy-dependent occupation is chosen independently of these Fermi
functions. The selected extracted-particle current is already positive;
no orientation reversal is applied.

\section{Total entropy, circulation and the integrated certificate}

The sample is circulant in its lead indices. Let
$u=\mathcal T_{01}$, $v=\mathcal T_{02}$,
$d_{01}=f_0-f_1>0$ and $d_{12}=f_1-f_2>0$. The two circulation sectors
contribute respectively the current-density vectors
\begin{equation}
 u(d_{01},d_{12},-d_{01}-d_{12}),\qquad
 v(d_{01}+d_{12},-d_{01},-d_{12}).
 \label{eq:cycles}
\end{equation}
Their sum conserves particle number. At the common temperature, conservation
of total energy and particle number converts the heat definition
\eqref{eq:entropy} into
\begin{equation}
 \frac{\sigma}{\kb}=\sum_a\frac{\mu_a-E_0}{\ee}J_a.
 \label{eq:commonT}
\end{equation}
In particular, define the dimensionless energy densities $\mathcal J$ and
$\mathcal E$ by
\begin{align}
 \mathcal J&=u d_{01}+v(d_{01}+d_{12}),\label{eq:Jdensity}\\
 \mathcal E&=u(25d_{01}+21d_{12})+v(4d_{01}+25d_{12}).
 \label{eq:Edensity}
\end{align}
Their energy integrals divided by $h$ give $J$ and $\sigma/\kb$.
Both are nonnegative, and the central circulating transmission makes them
strictly positive on a positive-width set. All three heats are included in
\eqref{eq:Edensity}.

For numerical evaluation use $q=(E/\ee-5000)/0.01$, with $\dd E=0.01\ee\dd q$.
Put $z=e^{0.01q}$, $e_0=e^{-1}$, $e_1=e^3$, $e_2=e^{24}$ and
$D_a=1+ze_a$. Then
\begin{equation}
 d_{01}=\frac{z(e_1-e_0)}{D_0D_1},\qquad
 d_{12}=\frac{z(e_2-e_1)}{D_1D_2},\qquad
 A_b=\frac{z[e_b(e^{-24}-1)+e_0(e^{24}-1)]}{D_0D_b}.
 \label{eq:densities}
\end{equation}
Equations~\eqref{eq:marginal}--\eqref{eq:M} give
$M(-24)=1+uA_1+vA_2$. The residual in \eqref{eq:Delta} is integrated
directly with density
\begin{equation}
 \mathcal D=u(-76d_{01}-84d_{12})
             +v(8d_{01}-76d_{12})-\log(1+uA_1+vA_2).
 \label{eq:residualdensity}
\end{equation}
This retains correlations between current and entropy. Subtracting loose
separate observable enclosures would yield a weaker result.

The full finite-Hamiltonian range construction is described in
Appendices~\ref{app:continuants} and \ref{app:arithmetic}. It covers
$q\in[-8192,16384]$, corresponding to
$E/\ee-5000\in[-81.92,163.84]$, using exactly 26,275 consecutive dyadic
panels. Range bounds apply to every energy in a panel, including resonances.
All 727 panels with zero-inclusive interval denominator bounds are retained
using the physical transmission bounds. Analytic Fermi tails enclose the
rest of the actual propagating band. The tail contribution to the absolute
residual uncertainty is below $1.12\times10^{-14}\ee/h$.

Table~\ref{tab:certificate} gives conservative outward decimal enclosures
obtained by exact-rational summation of the exported panel bounds and
directed tail estimates. They are intervals, not asserted significant digits
of a point estimate. More detailed exact endpoints are supplied with the
reproducibility data. In particular,
\begin{equation}
 \Delta\in[0.019447591396978867,\,0.024247460886202978]\,\frac{\ee}{h}.
 \label{eq:certifiedDelta}
\end{equation}
The enclosure width is below $0.005\ee/h$.

\begin{table}[t]
\centering
\caption{Outward enclosures for the fixed physical sample. All entries are
in units of $\ee/h$. The residual is enclosed directly, rather than by
subtracting the other rows.}
\label{tab:certificate}
\begin{tabular}{lrr}
\toprule
Quantity & Lower endpoint & Upper endpoint\\
\midrule
$\Delta$       & 0.019447 & 0.024248\\
$J$            & 0.013576 & 0.014625\\
$\sigma/\kb$   & 0.073446 & 0.079193\\
$\chi_0(-24)$  & 0.010467 & 0.012105\\
\bottomrule
\end{tabular}
\end{table}

\begin{proposition}
For the stationary sample and reservoirs in
Section~\ref{sec:model}, the curvature-matched assertion
\eqref{eq:assertion} is false. In the stated computer-assisted arithmetic
trust model,
\begin{equation}
 I_0(-J)-B_{\psi_0}(-J)>0.019\,\frac{\ee}{h}.
\end{equation}
\end{proposition}
\begin{proof}
The complete certificate proves $J,\sigma>0$ and
$\Delta>0.019\ee/h$. Appendix~\ref{app:psi} gives a positive finite
exact $\psi_0$, so $R_{\psi_0}>1$ and $-J$ is strictly inside the asserted
domain. Equation~\eqref{eq:Delta} proves the strict violation. One admissible
finite sample suffices to refute the universal quantifiers.
\end{proof}

\section{Verification, physical meaning and limitations}

The integral proof uses complete range coverage, rather than quadrature
convergence. An independently implemented check reconstructs the retarded
boundary and branches by outgoing Bloch modes and scalar continuants;
another uses forward branch elimination and a direct core inverse.
Explicit Slater minors agree with the original counting determinant.
The latter implementation checks 58 high-precision scattering/counting
points, 41 comparisons with full-panel bounds and five dense 387-site
resolvent solves. The maximum dense-solve scattering disagreement is below
$5.6\times10^{-15}$. Gauss orders 2, 4 and 8 on the frozen partition give a
diagnostic residual near $0.02382794\ee/h$, with order-4-to-8 difference
about $8.1\times10^{-9}\ee/h$. These agreements are diagnostics, not the
source of the strict integral error bound. A positive-temperature reciprocal
boxcar baseline also reproduces the published narrow-window current and
CGF limits. Parameters, primary inputs and arithmetic libraries are shared
between these implementations; code and scattering values are not reused
for the independent physical reconstruction.

The range verification has an explicit qualification. An intermediate
scalar absolute-value helper is nearest rounded and fails to be an outward
magnitude bound for 42 actual coefficients. The fixed calculation is
validated by an independent directed 500-digit reconstruction and an
exhaustive comparison of every implemented remainder upper endpoint:
52,550 checks, two polynomials on all 26,275 final panels, with zero
violations. Later directed arithmetic provides sufficient margin for these
actual remainders. This establishes the fixed certificate, not generic
outward correctness of the helper for other grids or parameters. The
failure records and complete closure checker are retained with the data.

The mechanism can be read directly from \eqref{eq:cycles}. Around the
central circulating transmission, particles flow from the high-chemical-
potential reservoir into the intermediate reservoir and then toward the
low-chemical-potential reservoir. The intermediate reservoir reduces the
total mean entropy cost per selected extracted particle. In contrast, a
negative selected count is sensitive to the occupation of the much lower-
chemical-potential reservoir. The large negative counting field accentuates
that rare event. Finite filtering and magnetic flux implement this
separation in a physical stationary sample. Leakage at the filter edges is
not negligible by assumption; it is included in the certificate.

Accordingly, matching the curvature of a variance inequality does not
determine the proposed full rare-current shape from actual mean entropy.
The result says nothing against the established reciprocal large-deviation
theorem or the nonreciprocal variance inequality. It does not resolve the
joint-current alternative, exclude other total-entropy-dependent shapes,
or assert that the sample has been fabricated. The filter order, width,
core coupling and lead band are fixed. There is no infinite-size,
zero-temperature, wide-band or perturbative limit in the proof.

The computation trusts Python exact integer/rational arithmetic,
mpmath 1.3.0 directed interval elementary functions and nearest-rounded
arbitrary-precision basic arithmetic with the explicit forward-error
envelope. It is a computer-assisted model counterexample, without
proof-assistant formalization. The numerical margin is conditional on that
documented computational trust base.

\paragraph{Computational research method.}
OpenAI Codex assisted literature discovery, model selection, the analytic
reduction, code generation, computational checking and interpretation.
The reported intervals are outputs of the deterministic scattering and
range calculations described above. Separate implementations and the
complete fixed-certificate remainder comparison supply the disclosed
validation; AI text or confidence is not a numerical error estimate.

\paragraph{Data availability.}
The accompanying local reproducibility package contains the defining
parameters, exact-model derivations, complete exported panel data,
calculation scripts, arithmetic-defect records and the exhaustive
remainder-check code and outputs. It includes instructions for scratch-copy
recomputation and exact resummation. No public repository or deposited-data
identifier is claimed for this draft.

\section*{Acknowledgements and AI disclosure}
OpenAI Codex was used for research ideation, literature discovery, analytic
reasoning, computational implementation and checks, and manuscript drafting.
The disclosed validation consists of deterministic calculations, complete
range/error checks and separate AI-assisted implementations. Human
scientific review, critical revision of the AI-assisted text and references,
and confirmation of authorship remain pending for this local draft.

\appendix
\section{An analytic lower bound for the exact constant}
\label{app:psi}

The numerator in \eqref{eq:psi} is the binary relative entropy
$\KL(y\Vert x)$. Jensen's inequality applied to the square-root likelihood
ratio gives
\begin{equation}
 \KL(y\Vert x)\ge-2\log C,\qquad
 C=\sqrt{xy}+\sqrt{(1-x)(1-y)}.
\end{equation}
Put $x=\sin^2\alpha$, $y=\sin^2\beta$ with
$0<\alpha,\beta<\pi/2$. First take $\alpha>\beta$, and define
$r=\sin(\alpha-\beta)$, $A=\sin(\alpha+\beta)$, so $0<r\le A\le1$.
Then
\begin{equation}
 C=\sqrt{1-r^2},\qquad x-y=Ar,\qquad
 x+y-x^2-y^2=\tfrac12[A^2(1-2r^2)+r^2].
\end{equation}
For fixed $r$, let
\begin{equation}
 F_r(A)=Ar\asinh\frac{2Ar}{A^2(1-2r^2)+r^2}.
\end{equation}
With $v=2Ar/[A^2(1-2r^2)+r^2]>0$ and $b=1-2r^2$, differentiation gives
\begin{equation}
 \frac{F_r'(A)}r=\asinh v+\frac{v}{\sqrt{1+v^2}}
 \frac{r^2-bA^2}{r^2+bA^2}
 \ge\asinh v-\frac{v}{\sqrt{1+v^2}}\ge0.
\end{equation}
The first inequality follows because the displayed rational factor is at
least $-1$; the second follows by differentiation from $v=0$.
Thus the denominator of \eqref{eq:psi} is at most
\begin{equation}
 F_r(1)=r\asinh\frac{2r}{1-r^2}=2r\atanh r.
\end{equation}
Also $g(r)=-\log(1-r^2)-r\atanh r$ satisfies $g(0)=0$ and
\begin{equation}
 g'(r)=\frac r{1-r^2}-\atanh r\ge0,
\end{equation}
since the derivative of the right-hand side is
$2r^2/(1-r^2)^2\ge0$ and its value at zero is zero. Combining the bounds,
\begin{equation}
 \KL(y\Vert x)\ge-\log(1-r^2)\ge r\atanh r
 \ge\tfrac12(x-y)\asinh\frac{x-y}{x+y-x^2-y^2}.
\end{equation}
The square-root likelihood bound holds in both relative-entropy directions;
$C$ and the denominator are symmetric. The argument therefore covers
$x<y$ as well as all continuous limits, proving $\psi_0\ge1/2$.
At $y=x+\delta$ the numerator and denominator both have leading term
$\delta^2/[2x(1-x)]$, so the diagonal limit is one and $\psi_0\le1$.

\section{Finite continuants and transmission ranges}
\label{app:continuants}

This appendix uses $0.01\ee$ as the energy scale, so the filter diagonal
variable is $q$. Define
\begin{align}
 \gamma_f&=\frac1{2\sin(\pi/256)},\quad g_c=20000000,\quad
 \ellambda^2=\gamma_fg_c,\nonumber\\
 b_\ell&=\frac1{4\sin[(2\ell+1)\pi/256]\sin[(2\ell-1)\pi/256]},
 \label{eq:qparams}
\end{align}
where $\ellambda$ denotes the dimensionless core-filter hopping.
The closed-chain continuant has recurrence
$P_{j+1}=qP_j-b_jP_{j-1}$, $P_0=1$, $P_1=q$.
Let $F$ be the 128-site determinant, $B$ the 127-site determinant with
the first site removed and $G$ the 126-site determinant with both endpoints
removed. The profile is reflection symmetric, so deleting the other endpoint
gives the same $B$.

For clarity define
\begin{equation}
 a_c=\frac{\Gamma_f/\ee}{\sqrt{5000\cdot35000}},\qquad
 \gamma(q)=\gamma_f\sqrt{(1+0.01q/5000)(1-0.01q/35000)}.
\end{equation}
The exact contact-modified determinants are
\begin{equation}
 D=F+(\ii\gamma-a_cq)B,\qquad
 C=qD-\ellambda^2[B+(\ii\gamma-a_cq)G].
\end{equation}
The core Schur complement is $(C/D)\mathbb I-\ii g_c(P-P^{\mathsf T})$.
Its three sector determinants are $C$, $C-\sqrt3g_cD$ and
$C+\sqrt3g_cD$. Put $H=\prod_{\ell=1}^{127}b_\ell$ and
\begin{equation}
 \mathcal A=\frac{4\gamma^2(\ellambda^2)^2H^2g_c^2}
 {|C|^2|C-\sqrt3g_cD|^2|C+\sqrt3g_cD|^2}.
\end{equation}
The endpoint cofactors and core adjugate give
\begin{equation}
 u=\mathcal A|C-\ii g_cD|^2,\qquad
 v=\mathcal A|C+\ii g_cD|^2.
 \label{eq:uv}
\end{equation}
These identities refer to the finite physical Hamiltonian. At $q=0$ they
give $u=0$, $v=1$, consistent with the current orientation above.

For the small reverse numerator the cancellation is performed before
interval evaluation. With $B_c=qB-\ellambda^2G$ and
$K=F+\gamma_f^2G$, the exact components are
\begin{align}
 \operatorname{Re}(C-\ii g_cD)
 &=qF+g_c(\gamma-\gamma_f)B-a_cqB_c,\\
 \operatorname{Im}(C-\ii g_cD)
 &=q(\gamma+g_ca_c)B
   -g_c(\gamma/\gamma_f)K+g_c(\gamma/\gamma_f-1)F.
\end{align}
All coefficients of $K$, including small nonleading ones, are retained.
No inferred zero, approximate filter identity or wide-band replacement is
used. If an interval denominator contains zero, the panel uses
$0\le u,v\le1$ from unitarity; it is not removed. The separate physical
bound $e^{-24}\le M(-24)\le e^{24}$ ensures a defined logarithm.

\section{Whole-panel arithmetic and the fixed-certificate closure}
\label{app:arithmetic}

The trigonometric hoppings and continuant coefficients are generated with
directed interval arithmetic at 250 decimal digits. Write
$F(q)=F_e(q^2)$ and $B(q)=qB_e(q^2)$. On a dyadic $q$ panel let
$x=q^2$, $x_c=m^2$ for its midpoint $m$, and
$r_x\ge\sup|x-x_c|$, $X\ge\sup|x|$. A degree-19 Taylor range for either
even polynomial $p(x)=\sum_kp_kx^k$ includes the remainder
\begin{equation}
 \mathcal R_p\le r_x^{20}\sum_{k\ge20}|p_k|\binom{k}{20}X^{k-20}.
 \label{eq:remainder}
\end{equation}
This bounds all omitted degrees throughout the panel.

The first 20 Taylor coefficients are evaluated by generalized Horner jets
using 336-bit nearest-rounded point arithmetic. The dyadic midpoint square
is exactly representable. Let $u_p=2^{1-336}$,
$N=5544$ and $\gamma_N=Nu_p/(1-Nu_p)$. The operation count exceeds all
basic operations in each signed or positive jet recurrence. If
$\widehat A$ is the computed positive recurrence for absolute coefficient
sums and $r_c$ is the directed relative coefficient uncertainty, the jet
envelope is
\begin{equation}
 \widehat{\mathrm{jet}}\ \pm\
 \frac{(\gamma_N+r_c)\widehat A}{1-\gamma_N}.
\end{equation}
It includes cancellation and rounding of the positive recurrence.
Subsequent Taylor evaluation, norms, divisions, square roots, exponentials,
logarithms, panel widths and sums use directed 85-digit interval operations.

For reproducibility it is essential to describe a defect in the original
remainder implementation. The scalar absolute-value conversion used to
initialize coefficient magnitudes is nearest rounded. Independent exact
endpoint comparisons show 44 failures to enclose the original coefficient
intervals, including 42 strict underestimates of the exact coefficient
magnitudes. This helper cannot be treated as an outward upper-bound
operation in general.

The closure is specific and complete. A separately transcribed checker
reconstructs coefficient intervals at 500 digits. It extracts absolute
endpoint bounds by changing their encoded sign, avoiding nearest-rounded
scalar absolute values. At every actual final panel, it evaluates a
500-digit directed reference bound for \eqref{eq:remainder} and compares
it to the implemented 85-digit remainder upper endpoint. Each implemented
upper endpoint is at least the reference upper endpoint, for both
polynomials on every one of the 26,275 panels. All 52,550 comparisons
complete with zero violations. Hence every remainder used in this fixed
certificate is adequate despite the defective intermediate initialization.
The supplemental code and records preserve both the failure and this
exhaustive closure; they do not certify a reusable helper for new inputs.

Every panel supplies a range for the full density. Its positive width times
that range encloses the integral, including structures narrower than the
panel. The refinement driver can retain valid parent intervals or replace
them by child ranges. The final consecutive dyadic endpoints and width
subtractions are checked as exact rationals, with total $q$ width 24576.
Each binary64 exported lower/upper endpoint is moved one representable
number outward. Exact-rational resummation covers every exported panel,
including all 727 denominator fallbacks, and adds the tails below once.

\section{Finite-band tails}
\label{app:tails}

Let $\delta=E/\ee-5000$. Below $-81.92$, every reservoir hole probability
is bounded by $e^{\delta+24}$; above $163.84$, every particle probability
is bounded by $e^{1-\delta}$. A nonzero extracted count requires at least
one hole in the low-energy tail or at least one particle in the high-energy
tail. Conservative dimensionless integrated event bounds are
\begin{equation}
 H_t=3e^{-81.92+24},\qquad P_t=3e^{1-163.84}.
\end{equation}
Throughout these tails the $e^{24}$-amplified event probability is below
$1/2$. From $|M-1|\le e^{24}\Pr(X\ne0)$ and
$|\log M|\le2|M-1|$, the following bounds hold in rate units $\ee/h$:
\begin{align}
 J_{\mathrm{tail}}&\le2(H_t+P_t),\\
 (\sigma/\kb)_{\mathrm{tail}}&\le28J_{\mathrm{tail}},\\
 |\chi(-24)|_{\mathrm{tail}}&\le2e^{24}(H_t+P_t),\\
 |\Delta|_{\mathrm{tail}}&\le24J_{\mathrm{tail}}
       +4(\sigma/\kb)_{\mathrm{tail}}+|\chi(-24)|_{\mathrm{tail}}.
\end{align}
Directed exponentials bound the last expression by
$1.1140\times10^{-14}\ee/h$. Extending the estimates beyond the actual
lead band only enlarges the bound. Together with the complete central
partition this covers all physical energies; no infinite-energy truncation
is left uncontrolled.

\bibliographystyle{unsrtnat}
\bibliography{references}
\end{document}
